\pdfcompresslevel=0
\pdfobjcompresslevel=0
\documentclass[11pt]{article}

\usepackage[a4paper,margin=1in]{geometry}
\usepackage[numbers,sort&compress]{natbib}
\usepackage{amsmath,amssymb,mathtools}
\usepackage{booktabs}
\usepackage{graphicx}
\usepackage{tabularx}
\usepackage{array}
\usepackage{multirow}
\usepackage{siunitx}
\usepackage{enumitem}
\usepackage{xcolor}
\usepackage{url}
\usepackage{microtype}
\usepackage{caption}
\usepackage{hyperref}

\sisetup{detect-all,group-separator={,},group-minimum-digits=4}
\hypersetup{
  colorlinks=true,
  citecolor=blue!55!black,
  linkcolor=blue!55!black,
  urlcolor=blue!55!black,
  pdftitle={State Sensitivity Is Not Grounding: Dual Replay for LLM Battery Scheduling},
  pdfauthor={Jiangwei Xue, Zhida Qin, Yuda Bi},
  pdfsubject={State grounding and raw-action feasibility in language-model battery scheduling},
  pdfkeywords={agent evaluation, battery scheduling, environmental grounding, large language models, physical feasibility, state intervention},
  pdfinfo={
    DOI={10.5281/zenodo.22707487},
    Copyright={Copyright 2026 The Authors},
    License={Creative Commons Attribution 4.0 International},
    LicenseURL={https://creativecommons.org/licenses/by/4.0/}
  }
}
\setlength{\emergencystretch}{2em}
\setlength{\textfloatsep}{10pt plus 2pt minus 2pt}
\setlength{\floatsep}{8pt plus 2pt minus 2pt}
\setlength{\intextsep}{8pt plus 2pt minus 2pt}
\setlength{\abovecaptionskip}{5pt}
\setlength{\belowcaptionskip}{2pt}
\captionsetup{font=small,labelfont=bf}
\setlist{nosep,leftmargin=*}
\raggedbottom

% Frozen result macros. Sources are the v1.7 offline evidence release and the
% paper-design package supplied with this draft.
\newcommand{\EIIBlocks}{360}
\newcommand{\EIIScenarios}{60}
\newcommand{\EIIPlannedCalls}{2160}
\newcommand{\EIIActualCalls}{2158}
\newcommand{\EIISuccessfulCalls}{2146}
\newcommand{\EIICS}{466.88}
\newcommand{\EIICSLow}{280.24}
\newcommand{\EIICSHigh}{671.80}
\newcommand{\EIINull}{207.13}
\newcommand{\EIINullLow}{88.75}
\newcommand{\EIINullHigh}{344.00}
\newcommand{\EIIAdjusted}{259.75}
\newcommand{\EIIAdjustedLow}{55.75}
\newcommand{\EIIAdjustedHigh}{464.80}
\newcommand{\EIIDeepAdjusted}{405.58}
\newcommand{\EIIDeepAdjustedLow}{120.99}
\newcommand{\EIIDeepAdjustedHigh}{704.76}
\newcommand{\EIIQwenAdjusted}{113.91}
\newcommand{\EIIQwenAdjustedLow}{-178.91}
\newcommand{\EIIQwenAdjustedHigh}{426.41}
\newcommand{\EIIPositive}{20}
\newcommand{\EIIZero}{34}
\newcommand{\EIINegative}{6}
\newcommand{\EIIParsedPairsCS}{234}
\newcommand{\EIIParsedPairsCC}{238}
\newcommand{\FOneBranches}{480}
\newcommand{\FZeroBranches}{120}
\newcommand{\DeepMeanED}{1173.33}
\newcommand{\DeepMedianED}{0}
\newcommand{\QwenMeanED}{345.45}
\newcommand{\QwenMedianED}{0}
\newcommand{\QwenThreeSevenMeanED}{1461.55}
\newcommand{\QwenThreeSevenMedianED}{154.40}
\newcommand{\PhysicsComparisons}{4620}
\newcommand{\PhysicsFixtures}{12}
\newcommand{\EIIProjectionTimeouts}{53}
\newcommand{\EIIProjectionRows}{1800}

\title{\vspace{-1.0em}State Sensitivity Is Not Grounding:\\
Dual Replay for LLM Battery Scheduling}
\author{%
Jiangwei Xue\\
\small Independent Researcher\\
\small \texttt{jw.xue@lucentia.company}
\and
Zhida Qin\\
\small Independent Researcher\\
\small \texttt{shin\_robotics@akane.waseda.jp}
\and
Yuda Bi\thanks{Corresponding author.}\\
\small TReNDS Center (GSU, GaTech, Emory)\\
\small \texttt{yudabi95@gmail.com}%
}
\date{%
\vspace{-0.6em}\small
DOI: \href{https://doi.org/10.5281/zenodo.22707487}{\textcolor{black}{10.5281/zenodo.22707487}}\\[0.45em]
\begin{minipage}{0.94\textwidth}
\centering\footnotesize
\textcopyright{} 2026 The Authors. This preprint is licensed under the
\href{https://creativecommons.org/licenses/by/4.0/}{\textcolor{black}{Creative Commons Attribution 4.0 International License (CC BY 4.0)}}.
\end{minipage}%
}

\begin{document}
\maketitle
\vspace{-1.2em}
\begin{abstract}
Battery schedules depend on the current state of charge (SOC). A model can react to an SOC change
yet produce an irrelevant or unsafe trajectory. We present an interventional audit for direct,
one-shot language-model scheduling. Forecasts, prices, constraints, events, and horizon remain
fixed while the represented SOC changes by 50, 75, or 100~kWh. Byte-identical repeated requests
estimate hosted-call variation. Each raw schedule is replayed from both the represented state and
the authoritative plant state. Across 60 challenge scenarios, two hosted Flash conditions
showed a mean canonical--stale distance of \EIICS{}~kWh versus \EIINull{}~kWh between duplicate
canonical calls. The null-adjusted difference was \EIIAdjusted{}~kWh
[\EIIAdjustedLow{}, \EIIAdjustedHigh{}] and was positive in \EIIPositive{}/60 scenarios. Under the
complete contract, DeepSeek mostly returned zero schedules and Qwen3.6 returned active schedules
with widespread SOC and export violations; neither produced an authoritative-feasible schedule in
480 branches. In a separate Qwen3.7-Plus tier, F1 yielded 6/240 carrier-consistent and 3/240
authoritative-feasible schedules, while the F0 payload subset yielded 1/60 and 0/60. Four schedules
passed replay under the represented state and failed under the authoritative state. A common
deterministic gate repaired every historical F1 branch. Matched zero and deterministic references
separated gate-created feasibility from proposal value. An independent checker reproduced all
\PhysicsComparisons{} physical verdicts with zero disagreement. The resulting
intervention-to-authority protocol measures state grounding in continuous physical decisions.
\end{abstract}

\paragraph{Keywords.} agent evaluation; battery scheduling; environmental grounding; large language models; physical feasibility; state intervention.

\section{Introduction}
\label{sec:introduction}

A grid-connected battery acts from a state that changes with every executed interval. State of
charge (SOC), usable capacity, converter limits, reserve commitments, export limits, forecasts,
and prices determine which schedules remain admissible. Model predictive control (MPC) makes this
dependence explicit by rebuilding the optimization problem from the current state and available
information \cite{parisio2014mpc,cedeno2022pvmpc,banfield2020empc}. Language models are now being
studied as energy interfaces, schedule generators, supervisory agents, and optimization surrogates
\cite{mongaillard2024power,michelon2025llmhems,cheng2025dispatch,elmakroum2026agentic,
raghavan2026costoptimal,haghighi2026surrogate}. Recent reviews and power-dispatch frameworks place
these systems on a spectrum from direct numerical generation to tool-mediated orchestration
\cite{zhang2026energysystems,zhao2026dispatchagent,zhao2026carbonagent}. This paper asks a narrower
question: when a numerical schedule changes after the represented physical state changes, what has
actually been measured?

A static feasibility test answers whether one output satisfies one instance. It does not reveal
whether the model used the state that made the instance different. Output variation alone also
leaves the question open. Hosted calls can vary under the same payload, and a large action change
can point away from the physically admissible region. In continuous physical scheduling, state
grounding therefore has several observable contracts: the action should respond beyond ordinary
repeated-call variation; the response should have a state-consistent direction; the schedule
should satisfy the physical contract under the state supplied to the model; and the same schedule
should remain valid when replayed from the authoritative plant state.

Several research lines motivate this measurement problem. Direct-optimization benchmarks test
whether language models can produce constraint-faithful schedules or optimization solutions under
deterministic verification \cite{tso2026constraintbench,jiang2026poweratlas,sharma2026schedbench}.
Agent and embodied-AI work studies action grounding, state maintenance, stale evidence, and
feedback-driven replanning \cite{huang2022zeroshot,ahn2022saycan,huang2023innermonologue,
yoneda2024statler,wang2025groundact,sheng2026envtrust,chao2026stale}. Repeated-trial benchmarks and
stochastic evaluation methods show that reliability cannot be inferred from one run
\cite{cui2026sameprompt,lior2025reliableeval,yao2024taubench}. Runtime-assurance research provides
separate mechanisms for enforcing physical constraints around learned or experimental controllers
\cite{seto1998simplex,alshiekh2018shielding,wabersich2021safetyfilter}. The measurement gap lies at
their intersection: a continuous physical schedule needs a controlled state intervention, a
matched repeat null, and replay under both the represented and authoritative dynamics.

We study that problem in a photovoltaic (PV)--battery environment with 15-minute actions and a
24-hour horizon. The experimental program contains four non-pooled evidence tiers. E2 identifies
state response by changing only the physical SOC carrier and comparing canonical--stale distance
with a byte-identical canonical duplicate. F1 supplies the complete battery contract and two
repeated calls per carrier. F0 repeats a fixed subset of the same payloads after approximately
6.4 days. Qwen3.7-Plus is retained as a later sensitivity tier with separate denominators. Every
proposal is scored before clipping, projection, or repair.

The paper makes four contributions.
\begin{enumerate}[leftmargin=*,itemsep=2pt,topsep=4pt]
  \item We formalize an intervention-to-authority audit that separates state response,
  directional consistency, carrier-consistent feasibility, authoritative feasibility, and
  post-gate proposal value.
  \item We combine a controlled 50/75/100-kWh SOC intervention with a byte-identical repeated-call
  null, identifying state-carrier effects without assuming deterministic hosted inference.
  \item We introduce dual replay for physical schedules. The same action is evaluated from the
  model-facing carrier state and from the authoritative state, exposing authority-transfer
  failures that a single replay cannot detect.
  \item We connect raw-action evidence to a matched execution-gate analysis using zero, heuristic,
  stale-optimal, random, and same-information MPC references. This locates the source of post-gate
  feasibility and economic performance.
\end{enumerate}

To our knowledge, prior work on direct scheduling, state grounding, and constraint-faithful agent
evaluation has not combined three elements in one physical scheduling protocol: a controlled
intervention on a decision-relevant continuous state, an identical-payload repeat null, and
deterministic replay under represented and authoritative states
\cite{jiang2026poweratlas,sharma2026schedbench,wang2025groundact,sheng2026envtrust,chao2026stale}.
The protocol covers open-loop schedule proposals under frozen hosted conditions. Tool-assisted and
closed-loop energy-management systems define related but different execution architectures.

\section{Related work}
\label{sec:related}

\subsection{Direct language-model optimization and energy scheduling}

Language models have been used in power scheduling as intent translators, direct dispatch agents,
residential schedulers, and surrogates for deterministic optimizers. Mongaillard et al. convert
user requests into optimization and allocation decisions, while Cheng et al. study an LLM for
advanced power dispatch \cite{mongaillard2024power,cheng2025dispatch}. Residential and building
studies connect language models with home-energy scheduling, MPC, or MILP-derived supervision
\cite{michelon2025llmhems,elmakroum2026agentic,raghavan2026costoptimal,haghighi2026surrogate}.
Other architectures reserve numerical optimization and verification for deterministic modules and
use the language model for orchestration or structured parameterization
\cite{su2026microgrid,zhao2026carbonagent}. A recent energy-systems review surveys this broader
landscape \cite{zhang2026energysystems}, and a power-dispatch-agent framework emphasizes planning,
memory, tool invocation, and action as separate capabilities \cite{zhao2026dispatchagent}.

Benchmark work increasingly isolates numerical and constraint faithfulness. ConstraintBench asks
models to return structured solutions to fully specified optimization problems without solver
access and checks every constraint against deterministic references \cite{tso2026constraintbench}.
PowerAtlas evaluates physical feasibility and economic outcomes of power schedules over
oracle-labeled instances \cite{jiang2026poweratlas}. SCHEDBench uses matched scheduling instances
to test whether semantically equivalent surface forms induce above-noise changes in hard-constraint
violations \cite{sharma2026schedbench}. PlanBench provides a broader planning benchmark grounded in
automated-planning domains \cite{valmeekam2023planbench}. Electricity-domain benchmarks extend the
scope to factuality, stability, live infrastructure, steady-state studies, and dynamic studies
\cite{zhou2024elecbench,curcio2026energyagentbench,mylonas2026poweragentss,
zhang2026poweragentdyn}. These studies establish feasibility and constraint faithfulness as central
evaluation targets. Our protocol focuses on the state dependence of one direct numerical schedule
and whether that dependence survives authoritative replay.

\subsection{Environmental state and action grounding}

Embodied-agent research provides a long line of work on mapping language to admissible actions and
updating plans as the environment changes. Zero-shot language-model planners often require explicit
grounding to executable actions \cite{huang2022zeroshot}; SayCan combines language scores with
robot affordances \cite{ahn2022saycan}; ProgPrompt and Code as Policies constrain or compile
language-model outputs into executable robot programs \cite{singh2023progprompt,
liang2023codepolicies}. ReAct and Inner Monologue interleave reasoning with observations or
feedback, allowing plans to update after interaction \cite{yao2023react,huang2023innermonologue}.
ALFWorld similarly aligns abstract text environments with embodied tasks
\cite{shridhar2021alfworld}.

State representation itself has also become an explicit object of study. Statler maintains a
structured world state across long-horizon embodied reasoning \cite{yoneda2024statler}. GroundAct
evaluates whether actions respect environmental prerequisites, while EnvTrustBench and STALE test
incorrect, stale, or superseded information \cite{wang2025groundact,sheng2026envtrust,
chao2026stale}. CausalWorld shows how controlled interventions on environment variables can support
causal and transfer evaluation in robotics \cite{ahmed2021causalworld}. Our setting adds continuous
battery dynamics and a counterfactual replay operation: one fixed trajectory can be evaluated from
two initial SOC values with identical forecasts, limits, and event information. The intervention
is graded in physical units, which supports direction and dose analyses.

\subsection{Repeated trials, verification, and execution gates}

Agent evaluation has moved beyond single-run success. AgentBench evaluates decision-making across
interactive environments, while AgentBoard exposes intermediate progress rather than only final
success \cite{liu2024agentbench,ma2024agentboard}. $\tau$-bench measures reliability over repeated
trials through pass$^k$ \cite{yao2024taubench}; Agent-SafetyBench evaluates safety failures across
interactive agent environments \cite{zhang2024agentsafety}. Repeated-execution studies and
stochastic evaluation methods likewise recommend distributions or repeat-aware estimators
\cite{cui2026sameprompt,lior2025reliableeval,liang2023helm}. We use a byte-identical canonical
request as an experimental null inside each E2 block, so state-carrier differences are compared
with variation observed under the same scheduling payload.

External verification and runtime assurance form a separate lineage. LLM-Modulo combines language
models with model-based verifiers \cite{kambhampati2024llmmodulo}. Simplex architectures separate
an experimental controller from a trusted safety controller \cite{seto1998simplex}; shielding and
predictive safety filters enforce state and input constraints around learned policies
\cite{alshiekh2018shielding,wabersich2021safetyfilter}. In energy control, MPC and economic MPC offer
deterministic references for constrained battery operation \cite{parisio2014mpc,banfield2020empc},
while battery scheduling can explicitly trade market value against degradation
\cite{sarker2017battery}. The execution gate in this study belongs to this runtime-assurance
tradition. We use it to ask how much of the final feasible trajectory is supplied by the proposal
and how much is created by the gate.

\section{Physical problem and grounding contracts}
\label{sec:formalization}

\subsection{PV--battery dynamics}

Each scenario $i$ contains an exogenous suffix
\begin{equation}
 z_i=\{\ell_{it},p^{\mathrm{pv}}_{it},\lambda_{it},\theta_{it}\}_{t=0}^{H_i-1},
 \label{eq:exogenous}
\end{equation}
where $\ell_{it}$ is load, $p^{\mathrm{pv}}_{it}$ is the PV forecast, $\lambda_{it}$ is energy
price, and $\theta_{it}$ collects time-varying battery and grid limits. The state includes SOC
$e_{it}$, usable capacity $E_{it}$, reserve $E^{\mathrm{res}}_{it}$, charge and discharge limits,
and the export limit.

The model returns a signed battery schedule
$u_i=(u_{i0},\ldots,u_{i,H_i-1})$, with positive values for charging. Let
$[a]_+=\max(a,0)$. For $\Delta t=0.25$~h, the deterministic SOC update is
\begin{equation}
 e_{i,t+1}=e_{it}+\eta_{\mathrm{ch}}[u_{it}]_+\Delta t
 -\frac{[-u_{it}]_+}{\eta_{\mathrm{dis}}}\Delta t.
 \label{eq:soc}
\end{equation}
Grid exchange, positive for import, is
\begin{equation}
 g_{it}=\ell_{it}-p^{\mathrm{pv}}_{it}+u_{it}.
 \label{eq:grid}
\end{equation}
A raw schedule is engineering-feasible when every interval satisfies the signed power bounds,
reserve and usable-capacity bounds on $e_{i,t+1}$, and the export condition
$g_{it}\ge -P^{\mathrm{exp}}_{it}$, while the terminal state satisfies
\begin{equation}
 |e_{i,H_i}-325|\le 1\ \mathrm{kWh}.
 \label{eq:terminal}
\end{equation}
The $\pm1$~kWh terminal tolerance defines the raw-action endpoint. The deterministic execution gate
uses an exact terminal equality. We retain these contracts as separate fields.

Let
\begin{equation}
 \Phi_i^{\epsilon}(u;x,z)\in\{0,1\}
 \label{eq:feasibility-map}
\end{equation}
denote deterministic replay from initial state $x$ with terminal tolerance $\epsilon$, and let
$\mathbf v_i^{\epsilon}(u;x,z)$ contain the associated power, SOC, export, and terminal violation
magnitudes. Binary feasibility and violation severity are both reported.

For post-gate economic comparisons, the physical cost is
\begin{equation}
 J_i(u)=\sum_{t=0}^{H_i-1}\left[
 \lambda_{it}g_{it}\frac{\Delta t}{1000}
 +c_{\mathrm{deg}}|u_{it}|\Delta t\right].
 \label{eq:cost}
\end{equation}
The prompt supplied prices and the complete physical contract. It did not instruct the model to
minimize Eq.~\eqref{eq:cost}. Cost and the same-information MPC are evaluation references applied
after the proposal is formed.

\subsection{Represented and authoritative state}

Let $x_i^A$ be the authoritative event-time state reconstructed from the executed prefix. The
canonical carrier contains $x_i^C=x_i^A$. The controlled stale carrier contains
\begin{equation}
 x_i^S=\mathcal I_{\delta_i}(x_i^A),
 \qquad |\delta_i|\in\{50,75,100\}\ \mathrm{kWh},
 \label{eq:intervention}
\end{equation}
with all non-carrier content held fixed. We use $x_i^R$ for the state represented to the model in
a given branch. Canonical branches have $x_i^R=x_i^A$; stale branches have $x_i^R=x_i^S$.

For model condition $m$, carrier $b\in\{C,S\}$, and repeated call $r$, the proposal is treated as
a hosted random output
\begin{equation}
 U_{imbr}\sim\Pi_m(\cdot\mid z_i,c(x_i^b)),
 \label{eq:model-output}
\end{equation}
where $c(\cdot)$ is the frozen state-carrier encoding.

\subsection{Five observable contracts}

We organize the audit through five observable contracts.

\paragraph{1. Response beyond repeat variation.}
For two schedules, define the action-energy distance
\begin{equation}
 d_i(u,v)=\Delta t\sum_{t=0}^{H_i-1}|u_t-v_t|.
 \label{eq:distance}
\end{equation}
E2 contains a byte-identical canonical duplicate $C2$. After reducing three repetitions to a
median within each scenario--model cell, the null-adjusted carrier response is
\begin{equation}
 \tau_{im}=D^{CS}_{im}-D^{CC2}_{im}.
 \label{eq:null-adjusted}
\end{equation}
A positive $\tau_{im}$ means that the carrier contrast exceeds the observed duplicate-call
contrast in that cell. For cross-scenario interpretation we also report
\begin{equation}
 \bar d_i(u,v)=\frac{d_i(u,v)}{H_i\Delta t},
 \qquad
 d_i^{\delta}(u,v)=\frac{d_i(u,v)}{|\delta_i|},
 \label{eq:normalized-distance}
\end{equation}
which express mean absolute power difference and action distance per kilowatt-hour of SOC
intervention.

F1 contains two calls per carrier. Its finite-sample two-carrier energy statistic is
\begin{align}
 B_{im}&=\frac{1}{4}\sum_{j=1}^{2}\sum_{k=1}^{2}d_i(C_j,S_k),\\
 W^C_{im}&=d_i(C_1,C_2),\qquad W^S_{im}=d_i(S_1,S_2),\\
 \mathrm{ED}_{im}&=2B_{im}-W^C_{im}-W^S_{im}.
 \label{eq:ed}
\end{align}
With two calls per group, $\mathrm{ED}_{im}$ can be negative. We report its distribution,
scenario support, and robust summaries alongside its mean.

\paragraph{2. Directional consistency.}
Let $\bar U_i^C$ and $\bar U_i^S$ be the two-call carrier means in F1, and let
$\delta_i=e_i^C-e_i^S$. A coarse directional score is
\begin{equation}
 R_i^{\mathrm{dir}}=-\operatorname{sgn}(\delta_i)\,
 \Delta t\sum_t(\bar U_{it}^C-\bar U_{it}^S).
 \label{eq:direction}
\end{equation}
Positive values indicate that a lower represented SOC leads to more net charging or less net
discharging. This score measures the sign of aggregate energy response; it does not encode the
full time-varying feasible set.

\paragraph{3. Carrier-consistent feasibility.}
The proposal is replayed from the model-facing state:
\begin{equation}
 F^{\mathrm{carrier}}_{imbr}=\Phi_i^1(U_{imbr};x_i^R,z_i).
 \label{eq:carrier-feasible}
\end{equation}
This endpoint asks whether the schedule is valid in the state represented to the model.

\paragraph{4. Authoritative feasibility.}
The same proposal is replayed from the plant state:
\begin{equation}
 F^{\mathrm{auth}}_{imbr}=\Phi_i^1(U_{imbr};x_i^A,z_i).
 \label{eq:auth-feasible}
\end{equation}
An authority-transfer failure is
\begin{equation}
 L^{\mathrm{auth}}_{imbr}=\mathbf{1}\{F^{\mathrm{carrier}}_{imbr}=1,
 F^{\mathrm{auth}}_{imbr}=0\}.
 \label{eq:authority-loss}
\end{equation}
We use $F^{\mathrm{auth}}$ as the safety-facing raw endpoint.

\paragraph{5. Proposal value after a common gate.}
Let $G_i(u)$ solve a deterministic feasibility-restoration problem from authoritative state. The
gate may alter battery power and, in the expanded formulation, use curtailment and mode variables.
Its post-gate regret is
\begin{equation}
 \operatorname{Regret}_i(G_i(u))=J_i(G_i(u))-J_i(u_i^{\star}),
 \label{eq:regret}
\end{equation}
where $u_i^{\star}$ is the same-information MPC reference. Matched proposal value against baseline
$b$ is measured by
\begin{equation}
 \Delta\operatorname{Regret}_{m,b}
 =\operatorname{Regret}_i(G_i(U_m))-\operatorname{Regret}_i(G_i(U_b)).
 \label{eq:proposal-value}
\end{equation}
A feasible $G_i(U_m)$ establishes feasibility of the proposal--gate system. Eq.~\eqref{eq:proposal-value}
locates any incremental economic value supplied by the proposal.

\section{Experimental design}
\label{sec:method}

\subsection{Public-data substrate and plant}

We assemble the deterministic environment from public 2024 time series at 15-minute resolution. PV input
uses NSRDB GOES CONUS v4 data for Los Angeles, San Diego, Fresno, and Sacramento
\cite{sengupta2018nsrdb}; \texttt{pvlib} converts irradiance and weather to AC power
\cite{anderson2023pvlib}. Load follows the OEDI ComStock LargeOffice profile
\cite{wilson2021eulp,parker2023comstock}. Price is CAISO day-ahead locational marginal price at
TH\_NP15\_GEN-APND \cite{caiso2017oasis}. The synchronized series define a controlled benchmark
substrate; they are not measurements from a co-located plant.

The simulated site combines a 400-kW PV array, a commercial load scaled to an 800-kW peak, and a
500-kWh battery. The battery accepts up to 250~kW in either direction. Both one-way efficiencies
are 0.95; the default reserve and export limits are 75~kWh and 150~kW. The battery-throughput
coefficient in Eq.~\eqref{eq:cost} is
0.008~USD/kWh; degradation-aware battery optimization motivates including throughput cost \cite{sarker2017battery}. A day contains 96 intervals. Each scenario contains one visible event from five
families: forecast revision, battery-capacity derating, power-limit derating, export-limit update,
and reserve-commitment update.

\subsection{Challenge-set construction}

The candidate generator produced 280 E2 scenarios. Provider-free admission required a feasible
canonical deterministic suffix, an infeasible zero suffix, at least 25~kWh difference between
canonical and stale deterministic plans, and a material consequence such as stale-plan
infeasibility, terminal error above 10~kWh, or a first-two-hour action change of at least 20~kW.
Two hundred five candidates passed admission, and 60 were selected with 12 scenarios per event
family. The set is designed to expose state-dependent failures. Its denominators describe this
challenge set. The authoritative SOC is lower than the stale carrier in 37 scenarios and higher in
23. Signed gaps are $-100$, $-75$, $-50$, $+50$, and $+100$~kWh.

\subsection{Evidence tiers}

Table~\ref{tab:tiers} summarizes the frozen tiers. They share the physical substrate and replay
code but answer different questions. Their results are reported separately.

\begin{table}[!htbp]
\centering
\caption{Experimental tiers and inferential roles. Calls are logical hosted requests.}
\label{tab:tiers}
\small
\begin{tabularx}{\textwidth}{@{}lrrrX@{}}
\toprule
Tier & Scenarios & Model conditions & Calls & Primary role \\
\midrule
E2 & 60 & 2 & 2,160 planned & Carrier response against a byte-identical null; three repetitions \\
F1 & 60 & 2 & 480 & Complete physical contract; two canonical and two stale calls \\
F0 & 15 & 2 & 120 & Byte-identical temporal rerun of an F1 subset after 6.44 days \\
Qwen3.7 F1 & 60 & 1 & 240 & Separate-model sensitivity tier under the F1 contract \\
Qwen3.7 F0 payloads & 15 & 1 & 60 & Separate sensitivity on the frozen F0 payload subset \\
\bottomrule
\end{tabularx}
\end{table}

\paragraph{E2 carrier identification.}
Each of the 60 scenarios is crossed with DeepSeek-V4-Flash and
Qwen3.6-Flash-2026-04-16 and run three times, producing \EIIBlocks{} blocks. Each block begins with
one shared prefix and then forks into five Stage-2 branches: canonical state (C), controlled stale
physical state (S), metadata control (K), model-maintained carrier (M), and a byte-identical
canonical duplicate (C2). The planned run contained \EIIPlannedCalls{} calls;
\EIIActualCalls{} were attempted and
\EIISuccessfulCalls{} returned successfully, with no protocol exclusions. C--S is the physical
intervention contrast and C--C2 is the repeated-call null used in the main estimator.

\paragraph{F1 complete-contract tier.}
F1 uses the same 60 scenario identities. Each scenario--model block contains C1, C2, S1, and S2,
for \FOneBranches{} calls. The prompt gives the signed power convention, SOC dynamics, efficiencies,
interval-specific charge/discharge, reserve, capacity and export limits, the full load/PV/price
suffix, and the 325-kWh terminal target with a 1-kWh engineering tolerance. Output is a sparse
JSON segment list that expands deterministically to the dense action vector. Temperature is zero,
provider reasoning is disabled where available, and no tools, solver, search, code execution, or
fallback route is permitted.

\paragraph{F0 temporal rerun.}
F0 resubmits 120 byte-identical payloads from 15 scenarios per historical model condition. The
first F0 request starts 154.65 hours after the last F1 request. This tier measures paired temporal
reproducibility across the two saved service windows.

\paragraph{Qwen3.7-Plus sensitivity.}
A later tier applies the same frozen F1 and F0 payload sets to Qwen3.7-Plus. It contains 240 F1
branches and 60 F0-subset branches. These rows are never pooled with the two historical Flash
conditions.

\subsection{Reduction and uncertainty}

Within each scenario--model cell, E2 takes the median of three repetition-level observations.
Model estimates average the resulting 60 cell medians. For the pooled estimate, the two model-cell
medians are averaged within scenario before averaging over the 60 scenarios. Parser failure
receives zero distance under the registered full-denominator rule; joint parseability is reported
beside the estimate.
Uncertainty uses 10,000 scenario-cluster resamples with seed 20260716. The resulting intervals
measure stability to the composition of the frozen challenge set.

F1 treats the scenario--model block as the independent unit. We report means, medians, positive/
zero/negative support, 10\% and 20\% trimmed means, and deletion of the largest one, two, and five
scenario effects. A random carrier-label pairing test provides an additional finite-sample null
for Eq.~\eqref{eq:ed}. Dose and direction are stratified by signed SOC gap and model.

\subsection{Replay, gate, and evidence validation}

Every raw response, parsed action, state carrier, replay, and gate trajectory is hash linked. Raw
replay is performed before any repair. The F1 battery-only gate seeks the minimum same-action-space
correction under exact terminal equality; the expanded gate can also use curtailment and mode
variables. Global distance optimality is certified for 409/480 historical F1 branches. The
remaining 71 branches retain a feasible mode-library upper bound on correction distance. The same
gate is applied to language-model, zero, heuristic, random static-range, stale-optimal, and
direct-MPC proposals.

A second physics checker independently implements SOC transition, power limits, reserve, usable
capacity, export, curtailment, terminal semantics, and cost. It imports no simulator, scorer, or
gate functions from the primary implementation. The release compares \PhysicsComparisons{} saved
raw, gated, and deterministic-witness trajectories and includes \PhysicsFixtures{} hand-specified
boundary fixtures.

\section{Results}
\label{sec:results}

\subsection{The SOC carrier changes some schedules beyond the repeat null}

Table~\ref{tab:e2-core} reports the E2 identification result. The mean canonical--stale distance
was \EIICS{}~kWh [\EIICSLow{}, \EIICSHigh{}]. Byte-identical canonical calls differed by
\EIINull{}~kWh [\EIINullLow{}, \EIINullHigh{}]. Their null-adjusted contrast was
\EIIAdjusted{}~kWh [\EIIAdjustedLow{}, \EIIAdjustedHigh{}]. Joint parsing succeeded for C/S in
\EIIParsedPairsCS{}/\EIIBlocks{} blocks and for C/C2 in
\EIIParsedPairsCC{}/\EIIBlocks{}.

% Evidence: manuscript/generated/e12_results_macros.tex in the paper-design package.
\begin{table}[!htbp]
\centering
\caption{E2 action-distance results. Values are kWh; brackets are scenario-cluster resampling
stability intervals.}
\label{tab:e2-core}
\small
\begin{tabular}{@{}lrr@{}}
\toprule
Estimate & Mean & Stability interval \\
\midrule
Canonical--stale, pooled & \EIICS & [\EIICSLow, \EIICSHigh] \\
Canonical--duplicate, pooled & \EIINull & [\EIINullLow, \EIINullHigh] \\
Null-adjusted, pooled & \EIIAdjusted & [\EIIAdjustedLow, \EIIAdjustedHigh] \\
Null-adjusted, DeepSeek-V4-Flash & \EIIDeepAdjusted & [\EIIDeepAdjustedLow, \EIIDeepAdjustedHigh] \\
Null-adjusted, Qwen3.6-Flash & \EIIQwenAdjusted & [\EIIQwenAdjustedLow, \EIIQwenAdjustedHigh] \\
\bottomrule
\end{tabular}
\end{table}

The model strata differ. DeepSeek's adjusted interval stays above zero; Qwen3.6's interval spans
zero. At scenario level, the pooled adjusted value is positive in \EIIPositive{}/60 cases, zero in
\EIIZero{}/60, and negative in \EIINegative{}/60. Capacity and power derating carry the largest
event-family means, while forecast and export updates are centered near zero. The aggregate effect reflects a subset of state-active cases rather
than a uniform response across the challenge set.

\subsection{Complete-contract responses are sparse and model-specific}

F1 changes the interface from E2's multi-branch experiment to a complete, direct schedule contract
with two calls per carrier. Table~\ref{tab:f1-phenotypes} shows three distinct output phenotypes.
DeepSeek returned 218/240 exact-zero schedules; Qwen3.6 returned active schedules in all 240
branches; Qwen3.7-Plus was active in 239/240 branches. Their response distributions differ as much
as their action levels.

% Evidence: expected_outputs/primary/f1_block_metrics.csv,
% expected_outputs/primary/f1_branch_metrics.csv, and Qwen37 counterparts.
\begin{table}[!htbp]
\centering
\caption{F1 response and authoritative-replay outcomes by model condition. Qwen3.7-Plus is a
separate sensitivity tier. ED support is positive/zero/negative over 60 scenario blocks.}
\label{tab:f1-phenotypes}
\small
\resizebox{\textwidth}{!}{%
\begin{tabular}{@{}lrrrrrrrr@{}}
\toprule
Model condition & Nontrivial & Mean ED & Median ED & ED support & Dynamic OK & Terminal viol. & $F^{\mathrm{carrier}}$ & $F^{\mathrm{auth}}$ \\
 & (/240) & (kWh) & (kWh) & ($+$/0/$-$) & (/240) & (/240) & (/240) & (/240) \\
\midrule
DeepSeek-V4-Flash & 22 & \DeepMeanED & \DeepMedianED & 8/52/0 & 223 & 240 & 0 & 0 \\
Qwen3.6-Flash & 240 & \QwenMeanED & \QwenMedianED & 3/56/1 & 1 & 240 & 0 & 0 \\
\midrule
Qwen3.7-Plus & 239 & \QwenThreeSevenMeanED & \QwenThreeSevenMedianED & 37/16/7 & 81 & 237 & 6 & 3 \\
\bottomrule
\end{tabular}}
\end{table}

DeepSeek's mean ED was \DeepMeanED{}~kWh, yet its median was zero and only 8/60 scenarios were
positive. Removing the five largest scenarios reduced the mean to 91.36~kWh; the 20\% trimmed mean
was zero. Qwen3.6 had a mean ED of \QwenMeanED{}~kWh, a zero median, and only three positive
scenarios. Removing its largest scenario moved the mean below zero. Random carrier-label pairing
placed the DeepSeek mean above its null distribution (upper-tail $p=0.00018$); the Qwen3.6 mean did
not ($p=0.13536$). Qwen3.7-Plus produced broader support, with a 154.40-kWh median and 37 positive
scenarios, while still retaining seven negative scenario statistics. After division by horizon
duration, mean ED is 51.56~kW for DeepSeek, 16.65~kW for Qwen3.6, and 64.86~kW for Qwen3.7-Plus;
the corresponding medians are 0, 0, and 6.80~kW.

The two historical Flash conditions fail in different ways. DeepSeek usually submits no action;
its dynamic constraints pass in 223/240 branches, but every branch misses the terminal target.
Qwen3.6 always acts, and 239/240 branches violate at least one dynamic constraint. SOC lower-bound
violations occur in 235 branches and export violations in 138. Both models finish with
0/240 authoritative-feasible schedules.

\subsection{Larger SOC interventions do not produce a monotone response}

Table~\ref{tab:dose} stratifies historical F1 by the absolute SOC gap. Every model--dose median is
zero. DeepSeek's mean ED rises from 50 to 75~kWh and then falls at 100~kWh. Qwen3.6 changes sign at
50~kWh and has its largest mean at 100~kWh. Directional support also remains sparse. The data do
not show a monotone dose--response relation over the three registered gaps. Event-family
composition differs across gap strata, so this table describes the registered challenge set rather
than an isolated causal dose effect.

% Evidence: expected_outputs/primary/f1_stratified_summary.csv.
\begin{table}[!htbp]
\centering
\caption{Historical F1 response by SOC intervention magnitude. Direction support counts positive,
zero, and negative coarse directional scores over 20 scenarios per cell.}
\label{tab:dose}
\small
\begin{tabular}{@{}lrrrrr@{}}
\toprule
Model & Gap (kWh) & Mean ED (kWh) & Median ED & ED support & Direction support \\
\midrule
\multirow{3}{*}{DeepSeek-V4-Flash}
 & 50  & 1137.50 & 0 & 2/18/0 & 1/18/1 \\
 & 75  & 1975.00 & 0 & 3/17/0 & 2/17/1 \\
 & 100 & 407.50  & 0 & 3/17/0 & 4/16/0 \\
\midrule
\multirow{3}{*}{Qwen3.6-Flash}
 & 50  & -155.25 & 0 & 0/19/1 & 2/15/3 \\
 & 75  & 126.50  & 0 & 1/19/0 & 1/13/6 \\
 & 100 & 1065.10 & 0 & 2/18/0 & 2/16/2 \\
\bottomrule
\end{tabular}
\end{table}

Across all gaps, DeepSeek's coarse directional score was positive in 7/60 scenarios, zero in 51,
and negative in two. Qwen3.6 was positive in five, zero in 44, and negative in 11. Positive action
distance alone provides little evidence that the response follows the sign of the SOC intervention.

\subsection{Carrier-consistent feasibility can fail under authoritative replay}

Table~\ref{tab:dual-replay} applies the two replay states to every complete-contract branch. The
historical Flash tiers remain infeasible under both definitions. Qwen3.7-Plus changes the picture:
6/240 F1 schedules are feasible from the state represented to the model, while 3/240 remain
feasible from the authoritative state. On the 60 F0-subset payloads, one schedule is
carrier-consistent and none is authoritative-feasible.

% Evidence: expected_outputs/semantics/FEASIBILITY_SEMANTICS.md and
% expected_outputs/qwen37/QWEN37_OFFLINE_STANDARDIZATION.json.
\begin{table}[!htbp]
\centering
\caption{Raw feasibility under the represented carrier state and authoritative state. Tiers are
kept separate.}
\label{tab:dual-replay}
\small
\begin{tabular}{@{}lrrr@{}}
\toprule
Tier & Branches & Carrier-consistent & Authoritative \\
\midrule
Historical F1, two Flash conditions & 480 & 0 & 0 \\
Historical F0 temporal rerun & 120 & 0 & 0 \\
Qwen3.7-Plus F1 & 240 & 6 & 3 \\
Qwen3.7-Plus F0 payload subset & 60 & 1 & 0 \\
\bottomrule
\end{tabular}
\end{table}

The four authority-transfer failures are all stale-carrier schedules. Three appear in Qwen3.7 F1
and one in its F0-subset tier. Their carrier-state terminal errors range from 0.03 to 0.71~kWh,
inside the engineering tolerance. Authoritative replay changes those errors to approximately
49.97, 50.19, 75.71, and 99.81~kWh. No other constraint creates these four mismatches. The fixed
action is internally valid for the represented initial SOC and invalid for the plant initial SOC.

\subsection{The temporal rerun preserves the feasibility conclusion}

All 120 historical F0 request payload hashes match their F1 counterparts. DeepSeek repeats 59/60
branch actions exactly and 14/15 four-branch blocks exactly. Qwen3.6 repeats 49/60 branch actions
and 9/15 blocks. Raw authoritative feasibility remains 0/120. This paired rerun shows that output
reproducibility is model-dependent even at temperature zero, while the physical endpoint is
unchanged across the two service windows.

\subsection{A common gate separates feasibility from proposal value}

The exact-terminal battery-only gate solves all 480 historical F1 branches. Post-gate feasibility
is 480/480, while raw authoritative feasibility is 0/480. The required corrections
are model-specific. DeepSeek's median correction is 131.58~kWh and Qwen3.6's is 5221.84~kWh. The
median post-gate regret is 3.12~USD for each model. The common median reflects the gate and scenario
set more than the raw action distribution.

Table~\ref{tab:gate-zero} reports matched differences against a zero proposal under the same gate.
DeepSeek has zero median differences in correction and regret; 218/240 correction differences and
224/240 regret differences are exactly zero. Qwen3.6 requires 5187.50~kWh more median correction
than zero. Its mean regret is 0.126~USD lower than zero, but the median is zero and the difference
is confined to 8/240 branches. The gate supplies feasibility in every branch. The raw proposal
supplies limited and comparator-dependent economic information.

% Evidence: OFFLINE_RECOMPUTE_SUMMARY.json, matched_llm_minus_baseline.
\begin{table}[!htbp]
\centering
\caption{Matched LLM-minus-zero differences under the battery-only gate. Support is
positive/zero/negative. Negative regret favors the LLM proposal. Intervals are
scenario-resampling stability intervals.}
\label{tab:gate-zero}
\small
\resizebox{\textwidth}{!}{%
\begin{tabular}{@{}lrrrrrr@{}}
\toprule
Model & Mean $\Delta D$ (kWh) & Median $\Delta D$ & $\Delta D$ support & Mean $\Delta$regret (USD) & Median & Regret support \\
\midrule
DeepSeek-V4-Flash & 337.58 [187.74, 510.46] & 0 & 17/218/5 & 0.253 [-0.330, 1.006] & 0 & 9/224/7 \\
Qwen3.6-Flash & 4633.93 [4489.48, 4774.67] & 5187.50 & 240/0/0 & -0.126 [-0.245, -0.035] & 0 & 0/232/8 \\
\bottomrule
\end{tabular}}
\end{table}

The expanded gate uses curtailment in 8/240 branches for each historical model. Most rescues alter
battery action and charge/discharge mode: 227/240 DeepSeek branches and 231/240 Qwen3.6 branches.
These counts show that repair can include control variables absent from the language-model action
space. Post-gate feasibility belongs to the full execution layer.

\subsection{Independent replay closes the physical evidence loop}

The independent checker agrees with the primary implementation on all \PhysicsComparisons{}
saved comparisons: 1,800 raw replay rows, 2,760 gate trajectories, and 60 deterministic witnesses.
All \PhysicsFixtures{} hand-specified fixtures pass, including power, SOC upper and lower bounds,
export, terminal tolerance, reserve, capacity derating, and boundary cases. The maximum numerical
difference is zero at the stored precision, and no Boolean feasibility verdict differs.

\section{Discussion}
\label{sec:discussion}

\subsection{State grounding is a sequence of contracts}

The experiments reveal distinct points of failure. E2 detects a carrier effect beyond an observed
repeat null, yet that effect is concentrated in a minority of scenarios and differs by model.
F1 shows that response magnitude can have a zero median and a large mean. Directional support is
sparser still. Qwen3.7-Plus provides the final separation: several schedules satisfy the complete
contract from the represented state, while four lose feasibility under authoritative replay.

These results support a contract view of state grounding. A physically grounded schedule should
show state-dependent behavior beyond repeat variation, align with the direction of the
intervention, and retain validity under the state that the plant will execute. Each contract has a
different evidence object. Response uses paired raw actions; direction uses signed intervention;
feasibility uses deterministic replay; authority transfer requires the same action under two
initial states. A single success rate would hide these distinctions.

\subsection{The repeat null changes the interpretation of sensitivity}

The E2 byte-identical distance is 44\% of the pooled canonical--stale distance. Treating identical
calls as deterministic would assign that entire difference to SOC. The within-block duplicate
provides a direct service-window reference. It also exposes model heterogeneity: DeepSeek's
adjusted interval stays above zero, while Qwen3.6's does not.

F1 reinforces the need for distributional reporting. DeepSeek's 1173.33-kWh mean ED is supported
by eight scenarios and a zero median. Qwen3.6 has one negative and three positive scenarios. These
patterns resemble thresholded or episodic state use more than a smooth controller response. The
50/75/100-kWh strata do not yield a monotone curve.

\subsection{Dual replay identifies a deployment-relevant failure}

Carrier-consistent replay measures internal validity under the state presented to the model.
Authoritative replay measures the action that would reach the plant. The four Qwen3.7 authority
losses are especially clear because the schedule has no nonterminal violation in either replay;
the initial SOC shift alone moves the terminal state outside tolerance. A conventional single
replay would classify these rows according to whichever initial state the evaluator happened to
choose.

The distinction applies beyond batteries. Any physical agent that receives an estimated,
cached, or model-maintained state can produce an action that is valid in that representation and
invalid in the execution state. Dual replay offers a general evaluation pattern when the dynamics
are deterministic enough to reconstruct both trajectories.

\subsection{Execution gates need matched attribution}

The F1 gate restores feasibility for every historical branch, including zero proposals and highly
active infeasible proposals. Its output combines the raw action, the exact-terminal contract, the
lexicographic objective, and additional variables such as curtailment and mode. The matched zero
comparison shows how little of the final result can be inferred from post-gate feasibility alone.
DeepSeek largely reproduces the zero baseline. Qwen3.6 occasionally changes cost by a small amount,
while requiring several megawatt-hours of additional trajectory correction.

This suggests a practical division of labor. Authoritative state belongs in a plant-side estimator
or digital twin. Deterministic optimization owns modeled feasibility. Language models can serve as
interfaces, intent translators, or proposal generators, as in solver-mediated and orchestration-focused energy architectures \cite{su2026microgrid,zhao2026carbonagent}. Their contribution should be measured
against simple proposals under the same execution layer. This comparison is more
informative than reporting a repaired success rate.

\subsection{Implications for benchmark design}

A physical-agent benchmark gains diagnostic power from four design choices. First, intervene on a
state variable whose physical consequence is known. Second, include identical-payload repeats in
the same service window. Third, preserve the represented state and the authoritative state as
separate replay inputs. Fourth, store raw and gated actions with stable row-level identifiers.
These choices support causal interpretation at the protocol level without requiring a model-based
judge.

The protocol also favors endpoint families over composite scores. Parseability, response,
direction, violation severity, authoritative feasibility, correction distance, and regret answer
different questions. Reporting them together preserves the path from model output to physical
execution.

\section{Limitations}
\label{sec:limitations}

The 60 scenarios form a purpose-built challenge set. Admission selects state-sensitive,
constraint-active cases, so the observed feasibility rates do not estimate ordinary PV--battery
operation. The candidate archive preserves selection metadata. Complete PV, load, and price vectors were
not retained for every rejected candidate, so a full counterfactual reselection cannot be
reconstructed.

The main hosted conditions are DeepSeek-V4-Flash and Qwen3.6-Flash under direct, one-shot,
structured-output prompts. Qwen3.7-Plus is a separate sensitivity tier. Reasoning tools, code
execution, solver calls, iterative correction, constrained decoding beyond JSON output, and
closed-loop replanning are outside the experiment. The results characterize schedule-proposal
reliability under these contracts.

E2 has three repetitions per scenario--model cell and a byte-identical null. F1 has two calls per
carrier, which makes Eq.~\eqref{eq:ed} variable and permits negative estimates. We report support,
trimmed summaries, deletion analyses, and a pairing null, but more repetitions would estimate the
full hosted-output distribution more precisely.

The prompt specifies feasibility and the terminal target. It does not specify Eq.~\eqref{eq:cost}
as the optimization objective. We exclude MPC trajectory distance as a measure of model
optimality.
Post-gate regret measures the execution system produced from each proposal.

The substrate synchronizes public PV, load, and price data from different sources and is not a
co-located field installation. The model proposes one open-loop suffix after one event. Sensor
error, battery ageing, network power flow, demand charges, hardware timing, receding-horizon
execution, and field protection remain untested.

The archived E2 projection population retains \EIIProjectionTimeouts{}/\EIIProjectionRows{}
timeouts and has status \texttt{PARTIAL\_WITH\_TIMEOUTS}. It is excluded from the complete-denominator
gate claims in this paper. The matched F1 gate population used in Section~\ref{sec:results} solves
all 480 branches.

\section{Conclusion}
\label{sec:conclusion}

This study evaluates whether direct language-model battery schedules use physical state in a form
that survives execution. A controlled SOC intervention changes some proposals beyond an
identical-payload repeat null, but the response is sparse, model-dependent, and non-monotone over
50, 75, and 100~kWh. The historical Flash conditions produce no authoritative-feasible schedule
under the complete contract. Qwen3.7-Plus produces a small set of carrier-consistent feasible
schedules, four of which fail when replay begins from the authoritative state. A deterministic
gate restores feasibility, while matched baselines show that the gate and the proposal contribute
different parts of the final result.

The central evaluation object is the path from state intervention to authoritative replay. This
path separates action sensitivity, directional use of state, physical validity in the represented
world, physical validity at the plant, and downstream proposal value. The resulting protocol is
applicable to other continuous physical agents whenever actions can be replayed under both
represented and authoritative dynamics.

\section*{Data and code availability}

Data and code supporting the findings are available in the accompanying reproducibility
repository at
\url{https://github.com/Jiangwei-Xue/paper-05-llm-battery-state-grounding}.
The repository provides the frozen tasks, prompts and request payloads, raw hosted-model
responses, parsed actions, deterministic dual replay, scores, branch comparisons, matched gate
baselines, claim-level hashes, analysis code, reference outputs, and a one-command offline
workflow that regenerates the reported analyses from saved provider records and validates the
VCR/H5/E5 evidence chain. The workflow does not repeat nondeterministic hosted-model calls.
The analysis code is released under the MIT License. Original project documentation and derived
result tables are licensed under CC BY 4.0. Third-party data, hosted-model records, and bundled
dependencies remain subject to their respective terms.

\section*{Competing interests}
The authors declare no competing interests.

\section*{Funding}
This research received no specific grant from any funding agency in the public, commercial, or
not-for-profit sectors.

\appendix

\section{Frozen schedule prompt contract}
\label{app:prompt}

The complete-contract prompt asks for one JSON object with a single \texttt{actions} field. Each
segment contains \texttt{start\_offset}, \texttt{end\_offset\_exclusive}, and \texttt{power\_kw}.
Segments must be sorted and non-overlapping; uncovered offsets expand to zero. Positive power
charges the battery and negative power discharges it. The prompt states Eq.~\eqref{eq:soc}, all
interval-specific power, reserve, capacity and export limits, Eq.~\eqref{eq:grid}, and the
terminal condition in Eq.~\eqref{eq:terminal}. It also states that raw actions are scored before
clipping, projection, or repair. The task JSON contains the carrier state and complete suffix time
series. Explanations and fields outside the schema are disallowed.

\section{Event-family heterogeneity in E2}
\label{app:event}

Table~\ref{tab:event} reports post-hoc event-family estimates. Each row contains 12 scenario
clusters. The intervals describe challenge-set resampling stability.

\begin{table}[!htbp]
\centering
\caption{E2 null-adjusted action distance by event family.}
\label{tab:event}
\small
\begin{tabular}{@{}lrrr@{}}
\toprule
Event family & Mean (kWh) & Stability interval & Support ($+$/0/$-$) \\
\midrule
Battery-capacity derating & 831.7 & [221.6, 1425.8] & 8/3/1 \\
Export-limit update & -67.5 & [-312.5, 156.7] & 2/8/2 \\
Forecast revision & -1.0 & [-550.0, 557.3] & 2/8/2 \\
Power-limit derating & 421.0 & [133.1, 734.4] & 5/6/1 \\
Reserve-commitment update & 114.6 & [0.0, 264.6] & 3/9/0 \\
\bottomrule
\end{tabular}
\end{table}

\section{Authority-transfer witnesses}
\label{app:witnesses}

Table~\ref{tab:witnesses} lists the four schedules with
$F^{\mathrm{carrier}}=1$ and $F^{\mathrm{auth}}=0$. All are Qwen3.7-Plus stale-carrier outputs.
The signed gap is canonical SOC minus stale SOC.

\begin{table}[!htbp]
\centering
\caption{Saved authority-transfer failures. Terminal errors are absolute kWh.}
\label{tab:witnesses}
\small
\resizebox{\textwidth}{!}{%
\begin{tabular}{@{}llllrrr@{}}
\toprule
Tier & Branch & Event & Location/season & Signed SOC gap & Carrier error & Authoritative error \\
\midrule
F1 & S1 & Forecast revision & Sacramento/winter & -75 & 0.71 & 75.71 \\
F1 & S1 & Export-limit update & Fresno/spring & +100 & 0.19 & 99.81 \\
F1 & S1 & Capacity derating & Sacramento/autumn & -50 & 0.19 & 50.19 \\
F0 subset & S2 & Export-limit update & San Diego/autumn & -50 & 0.03 & 49.97 \\
\bottomrule
\end{tabular}}
\end{table}

\section{Deterministic reference baselines}
\label{app:baselines}

Table~\ref{tab:baselines} gives aggregate battery-only gate outcomes for the five deterministic
proposal sources over the 60 F1 scenarios. These rows use the same authoritative state and exact
terminal gate. Direct MPC is the zero-regret reference by construction.

\begin{table}[!htbp]
\centering
\caption{Deterministic proposal baselines under the common battery-only gate.}
\label{tab:baselines}
\small
\resizebox{\textwidth}{!}{%
\begin{tabular}{@{}lrrrr@{}}
\toprule
Proposal source & Raw feasible (/60) & Median correction (kWh) & Median cost (USD) & Median regret (USD) \\
\midrule
Direct MPC & 60 & 0.00 & 337.38 & 0.00 \\
Simple heuristic & 60 & 0.00 & 354.47 & 5.15 \\
Stale-optimal plan & 0 & 71.25 & 337.38 & 0.00 \\
Zero schedule & 0 & 131.58 & 353.16 & 3.31 \\
Random static-range schedule & 0 & 368.50 & 361.49 & 27.23 \\
\bottomrule
\end{tabular}}
\end{table}

\begingroup
\small
\bibliographystyle{unsrtnat}
\bibliography{references}
\endgroup

\end{document}
